% ================================================================
%  第六章  财务预测与融资计划
%  写作建议：请专业人士制作，资金等关键数据前后一致。
%  核心聚焦：未来6~12个月需要多少钱？释放多少股权？用于什么？达成什么目标？
%  如项目尚未成熟，避免深度预测5年财务，重点说清楚当下融资逻辑。
% ================================================================

\chapter{财务预测与融资计划}

% ----------------------------------------------------------------
% 6.1 财务分析
%   · 历史财务数据（有则附，无则说明当前阶段）
%   · 未来3年收入预测与假设
%   · 成本结构分析（固定成本/变动成本）
%   · 盈利预测与现金流分析
%   · 投资计划与回报预期
% ----------------------------------------------------------------
\section{财务分析}

% ---- 6.1.1 历史数据与现状 ----
% TODO:
%   - 历史收入（如有）：
%     202X年：营业收入 XXX 万元，净利润 XXX 万元
%     202X年：营业收入 XXX 万元，净利润 XXX 万元
%   - 如尚未产生收入，说明当前阶段：
%     "项目目前处于产品研发阶段，已获得XXX验证，预计202X年开始商业化"
%   - 已投入资金：团队自投 XXX 万元，用于 XXX

% ---- 6.1.2 未来3年财务预测 ----
% TODO:
%   核心假设说明（非常重要！让评委知道数字从哪来）：
%   - 客单价假设：B端客户年均合同额 XXX 万元；C端用户月均消费 XXX 元
%   - 获客增长假设：第1年签约客户 XX 家/用户 XX 万，增速 XX%/年
%   - 成本结构假设：人力成本 XX%，服务器/原材料 XX%，市场费用 XX%
%
%   营收预测表（取消注释并填入数字）：
%   \begin{table}[htbp]
%     \centering
%     \caption{未来三年财务预测（万元）}
%     \begin{tabular}{lccc}
%       \toprule
%       科目 & 第1年（202X） & 第2年（202X） & 第3年（202X） \\
%       \midrule
%       营业收入 & & & \\
%       营业成本 & & & \\
%       毛利润  & & & \\
%       毛利率  & & & \\
%       运营费用（含研发+市场+管理） & & & \\
%       EBITDA  & & & \\
%       净利润  & & & \\
%       \bottomrule
%     \end{tabular}
%   \end{table}
%
%   现金流预测：
%   - 经营活动现金流：第1年 XXX 万元，第2年 XXX 万元
%   - 预计首次转正时间点：202X年X月（盈亏平衡月份）

% ---- 6.1.3 投资计划 ----
% TODO:
%   - 本次拟融资金额：XXX 万元
%   - 融资方式：股权融资 / 可转债 / 政府补贴+股权融资
%   - 出让股比：XX%，对应估值：XXX 万元
%   - 估值逻辑：基于【市盈率PE / 市销率PS / 对标估值法】，
%     对标企业 XXX 的估值倍数为 XX 倍，本项目类似规模下估值合理区间为 XXX 万元
%   - 抵押/担保条件：（如涉及银行贷款请说明）
%   - 投资收益预期：假设202X年收入达到XXX，以XX倍PS估值，届时退出估值约XXX万元
%   - 投后股权比例：
%       创始团队：XX%，投资方：XX%，期权池：XX%
%   - 财务报告编制频率：月报/季报/年报
%   - 投资人介入程度：董事会席位X个，日常经营自主，重大决策协商

% ---- 6.1.4 融资需求 ----
% TODO:
%   - 资金需求总额：XXX 万元（含本次融资XXX万 + 政府补贴XXX万 + 自有资金XXX万）
%   - 历史融资情况（如有）：
%       天使轮：202X年，XXX 万元，投资方 XXX，出让 XX%
%   - 资金用途明细表：
%   \begin{table}[htbp]
%     \centering
%     \caption{融资资金用途说明}
%     \begin{tabular}{lcc}
%       \toprule
%       用途类别 & 金额（万元） & 占比 \\
%       \midrule
%       产品研发（人力+设备） & & \\
%       市场推广与销售 & & \\
%       团队扩张（招聘） & & \\
%       日常运营（办公+行政） & & \\
%       预留流动资金 & & \\
%       \midrule
%       合计 & & 100\% \\
%       \bottomrule
%     \end{tabular}
%   \end{table}
%   - 融资方案：优先接受XXX类型投资人（产业投资/VC/战略合作方）

% ----------------------------------------------------------------
% 6.2 完成研发所需投入与盈亏平衡
%   · 研发总投入拆分（阶段化）
%   · 运营/市场/团队各项投入预算
%   · 盈亏平衡点计算与时间节点
%   · 项目里程碑与资金配置对应关系
% ----------------------------------------------------------------
\section{完成研发所需投入与盈亏平衡}
% TODO:
%   【分阶段研发投入】
%   - 第一阶段（202X.XX—202X.XX）：完成XXX，资金需求 XXX 万元
%       人力：XXX 万元（X人×X月×X元/月）
%       设备/软件：XXX 万元
%       外包/采购：XXX 万元
%   - 第二阶段（202X.XX—202X.XX）：完成XXX，资金需求 XXX 万元
%   - 第三阶段（202X.XX—202X.XX）：MVP上线/量产，资金需求 XXX 万元
%
%   【盈亏平衡计算】
%   - 固定成本（月）：XXX 万元（人力 + 场租 + 平台费等）
%   - 变动成本（每单位/每用户）：XXX 元
%   - 盈亏平衡销量/用户数：固定成本 ÷ (单价 − 单位变动成本) = XXX
%   - 预计达到盈亏平衡的时间节点：202X年X月（第X年第X季度）
%
%   【里程碑与资金配置时间表】
%   \begin{table}[htbp]
%     \centering
%     \caption{里程碑与资金配置}
%     \begin{tabular}{lll}
%       \toprule
%       时间节点 & 关键里程碑目标 & 资金需求 \\
%       \midrule
%       202X Q1 & & \\
%       202X Q2 & & \\
%       202X Q4 & & \\
%       202X Q4 & & \\
%       \bottomrule
%     \end{tabular}
%   \end{table}
